% -*- coding: utf-8 -*-
% !TEX program = xelatex
\documentclass[plain,noanswer,medmath]{../../template/jnuexam}
\usepackage{../../template/kysx}

\begin{document}


\examtitle{1991}{1}


\makepart{填空题}{本题满分15分，每小题3分}


\begin{problem}
设$\begin{cases}
  x = 1 + t^2, \\
  y = \cos t,
\end{cases}$ 则$\frac{\d^2y}{\dx^2} =$ \fillin{}．
\end{problem}


\begin{problem}
由方程$xyz + \sqrt{x^2 + y^2 + z^2} = \sqrt{2}$所确定的函数$z = z(x,y)$
在点$(1,0, - 1)$处的全微分$\dz =$ \fillin{}．
\end{problem}


\begin{problem}
已知两条直线的方程是$L_1:\frac{x - 1}{1} = \frac{y - 2}{0} = \frac{z - 3}{-1}$；
$L_2:\frac{x + 2}{2} = \frac{y - 1}{1} = \frac{z}{1}$，
则过$L_1$且平行于$L_2$的平面方程是 \fillin{}．
\end{problem}


\begin{problem}
已知当$x \to 0$时，$(1 + ax^2)^{\frac{1}{3}} - 1$与$\cos x - 1$是等价无穷小，
则常数$a =$ \fillin{}．
\end{problem}


\begin{problem}
设$4$阶方阵$A=\begin{pmatrix}
 5 & 2 & 0 &  0 \\
 2 & 1 & 0 &  0 \\
 0 & 0 & 1 & -2 \\
 0 & 0 & 1 &  1 \\
\end{pmatrix}$，则$A$的逆矩阵$A^{-1}=$ \fillin{}．
\end{problem}


\makepart{选择题}{本题满分15分，每小题3分}


\begin{problem}
曲线$y = \frac{1 + \e^{- x^2}}{1 - \e^{- x^2}}$	\pickout{}
\begin{abcd*}
\item 没有渐近线．					
\item 仅有水平渐近线．
\item 仅有铅直渐近线．				
\item 既有水平渐近线又有铅直渐近线．
\end{abcd*}
\end{problem}


\begin{problem}
若连续函数$f(x)$满足关系式$f(x) = \int_0^{2x} f\left(\frac{t}{2} \right) \dt + \ln 2$，
则$f(x)$等于\pickout{}
\begin{abcd}
\item $\e^x\ln 2$．		
\item $\e^{2x}\ln 2$．		
\item $\e^x + \ln 2$．		
\item $\e^{2x} + \ln 2$．
\end{abcd}
\end{problem}


\begin{problem}
已知级数$\sum_{n = 1}^{\infty}(- 1)^{n - 1} a_n = 2$，$\sum_{n = 1}^{\infty}a_{2n - 1} = 5$，
则级数$\sum_{n = 1}^{\infty}a_n$等于\pickout{}
\begin{abcd}
\item $3$．			
\item $7$．			
\item $8$．			
\item $9$．
\end{abcd}
\end{problem}


\begin{problem}
设$D$是$xOy$平面上以$(1,1)$，$(- 1,1)$和$(- 1, - 1)$为顶点的三角形区域，
$D_1$是$D$在第一象限的部分，则$\iint_D (xy + \cos x\sin y)\dx\dy$等于\pickout{}
\begin{abcd}
\item $2\iint_{D_1} \cos x\sin y\dx\dy$．
\item $2\iint_{D_1} xy\dx\dy$．
\item $4\iint_{D_1} (xy + \cos x\sin y)\dx\dy$．
\item $0$．
\end{abcd}
\end{problem}


\begin{problem}
设$n$阶方阵$A$、$B$、$C$满足关系式$ABC = E$，其中$E$是$n$阶单位阵，则必有\pickout{}
\begin{abcd}
\item $ACB = E$．    
\item $CBA = E$．    
\item $BAC = E$．    
\item $BCA = E$．
\end{abcd}
\end{problem}


\makepart{计算题}{本题满分15分，每小题5分}


\begin{problem}
求$\lim_{x \to 0^+} (\cos \sqrt{x})^{\frac{\pi}{x}}$．
\end{problem}


\begin{problem}
设$\vec{n}$是曲面$2x^2 + 3y^2 + z^2 = 6$在点$P(1,1,1)$处的指向外侧的法向量，求函数
$u = \frac{\sqrt{6x^2 + 8y^2}}{z}$在点$P$处沿方向$\vec{n}$的方向导数．
\end{problem}


\begin{problem}
求$\iiint_{\Omega}(x^2 + y^2 + z)\d V$，其中$\Omega$是由曲线$\begin{cases}
  y^2 = 2z, \\
  x = 0
\end{cases}$绕$z$轴旋转一周而成的曲面与平面$z = 4$所围成的立体．
\end{problem}


\makepart{}{本题满分6分}

\begin{problem}
在过点$O(0,0)$和$A(\pi ,0)$的曲线族$y = a\sin x(a > 0)$中，求一条曲线$L$，
使沿该曲线从$O$到$A$的积分$\int_L (1 + y^3)\dx + (2x + y)\dy$的值最小．
\end{problem}


\makepart{}{本题满分8分}

\begin{problem}
将函数$f(x) = 2 + |x|$ ( $- 1 \le x \le 1$)展开成以$2$为周期的傅立叶级数，并由此求级数
$\sum_{n = 1}^{\infty}\frac{1}{n^2}$的和．
\end{problem}


\makepart{}{本题满分7分}

\begin{problem}
设函数$f(x)$在$[0,1]$上连续，$(0,1)$内可导，且$3\int_{\frac{2}{3}}^1f(x)\dx = f(0)$，
证明在$(0,1)$内存在一点$c$，使$f'(c) = 0$．
\end{problem}


\makepart{}{本题满分8分}

\begin{problem}
已知$\alpha_1 = (1,0,2,3)$，$\alpha_2 = (1,1,3,5)$，$\alpha_3 = (1, - 1,a + 2,1)$，
$\alpha_4 = (1,2,4,a + 8)$，及$\beta = (1,1,b + 3,5)$．
\begin{enumerate}
\item $a$, $b$为何值时，$\beta$不能表示成$\alpha_1$, $\alpha_2$, $\alpha_3$, $\alpha_4$的线性组合？
\item $a$, $b$为何值时，$\beta$有$\alpha_1$, $\alpha_2$, $\alpha_3$, $\alpha_4$的唯一的线性表示式？
      并写出该表示式．
\end{enumerate}
\end{problem}


\makepart{}{本题满分6分}

\begin{problem}
设$A$为$n$阶正定阵，$E$是$n$阶单位阵，证明$A + E$的行列式大于$1$．
\end{problem}


\makepart{}{本题满分8分}

\begin{problem}
在上半平面求一条向上凹的曲线，其上任一点$P(x,y)$处的曲率等于此曲线在该点的法线段
$PQ$长度的倒数($Q$是法线与$x$轴的交点)，且曲线在点$(1,1)$处的切线与$x$轴平行．
\end{problem}


\makepart{填空题}{本题满分6分，每小题3分}


\begin{problem}
若随机变量$X$服从均值为$2$，方差为$\sigma^2$的正态分布，
且$P\{2 < X < 4\} = 0.3$，则$P\{X < 0\} =$ \fillin{}．
\end{problem}


\begin{problem}
随机地向半圆$0 < y < \sqrt{2ax - x^2}$（$a$为正常数）内掷一点，
点落在半圆内任何区域的概率与区域的面积成正比，
则原点和该点的连线与$x$轴的夹角小于$\frac{\pi}{4}$的概率为 \fillin{}．
\end{problem}


\makepart{}{本题满分6分}

\begin{problem}
设二维随机变量$(X,Y)$的概率密度为 
$$f(x,y) = \begin{cases}
 2\e^{-(x + 2y)}, & x > 0, y > 0; \\
 0, & \text{其他}.
\end{cases}$$
求随机变量$Z = X + 2Y$的分布函数．
\end{problem}


\end{document}
